\section{Introduction}
\label{sec:intro}
Submarine telecommunication cables constitute a critical infrastructure for global data transmission, making their accurate detection and localization essential for maintenance, inspection, and protection operations. However, identifying these cables in complex underwater environments remains a challenging task due to burial, environmental noise, and limited visibility.

Traditional sensing approaches rely primarily on acoustic and optical techniques. Acoustic systems, such as side-scan sonar, are effective for detecting exposed cables or those with visible seabed features, but their performance significantly degrades when cables are buried beneath sediments. Similarly, optical methods are constrained by poor underwater visibility, light attenuation, and interference from environmental artifacts, limiting their applicability in real-world conditions \cite{feng2024prediction,unal2024survey,leblond:hal-02304729}.

Magnetic sensing has become a widely used alternative \cite{oehler2024uuv}. Magnetometers detect submarine cables via magnetic anomalies from their ferromagnetic components or carried current, enabling localization even when the cables are buried. Existing cable and pipeline localization methods assume high SNR and use signal amplitude for detection \cite{Bharti2020,Yu2018,Xiang2016,li2024attitude}. Classical signal processing approaches for detection use multi-step pipelines—field separation, denoising, and inversion (e.g., Euler deconvolution \cite{Thompson1982EULDPHAN,reid2014avoidable})—that require careful parameter tuning, substantial human intervention, and can yield spurious solutions \cite{liu2024submarine}, limiting robustness and real-time use.

More recently, machine learning techniques have been introduced for magnetic anomaly detection and inversion. These data-driven approaches enable end-to-end mapping from raw measurements to cable positions and demonstrate improved robustness to noise and higher computational efficiency \cite{https://doi.org/10.1049/smt2.12084,HU2021106653,liu2024submarine}. However, their performance strongly depends on the availability and representativeness of training data, and their generalization capability may be limited when operating conditions differ from those encountered during training.

In parallel, orthonormal basis functions (OBF) have been widely used for the detection of magnetic dipole sources \cite{Otnes2007,liu2024adaptive}, providing a compact and physically meaningful description of the signal. These approaches offer advantages in terms of interpretability and computational efficiency, and remain relevant in modern detection frameworks \cite{https://doi.org/10.1049/rsn2.70091,fan2020adaptive,guo2024obf}. 
\highlight{
This framework was subsequently generalized to multipolar sources, first by Pepe et al.\ \cite{pepe2015etude} and later by Chenevas-Paule et al.\ \cite{chenevas2025multipolar}, where multipolar orthonormal basis functions (MOBF) were coined for detection, though these approaches address spatially bounded (compact) sources, enclosed within a finite Brillouin sphere \footnote{The smallest sphere containing the source}, which differs fundamentally from the cable geometry considered here.
}

In this context, the present study proposes the derivation of orthonormal basis functions specifically adapted to \highlight{infinitely long conductors}. By extending the OBF framework beyond the conventional \highlight{compact-source assumption}, the proposed methodology seeks to establish an efficient and physically grounded method for submarine telecommunications cable detection, thereby enabling the transfer and adaptation of OBF-based techniques originally developed for \highlight{multipoles} detection to the problem of cable detection.